% Generated from locale/tr/content/lambda-calculus/lambda-definability/arithmetical-functions.tex; do not hand-edit.


\olfileid[tr]{lam}{rep}{arf}
\olsection{\usetoken{S}{lambda definable} Aritmetik Fonksiyonlar}

\begin{prop}
  \ollabel{prop:succ-ld}
  Ardıl fonksiyonu $\Succ$ !!{lambda definable}.
\end{prop}

\begin{proof}
Ardıl fonksiyonunu !!{lambda define}s bir terim şudur:
\begin{align*}
  \fn{Succ} & \ident \lambd[a][\lambd[fx][f ({a} f x)]].
  \intertext{Uzlaşımlarımıza göre bu, şunun kısaltmasıdır:}
  \fn{Succ} & \ident \lambd[a][\lambd[f][\lambd[x][(f (({a} f) x))]]].
  \intertext{$\fn{Succ}$, sayı $a$'yı argüman olarak alan ve başka bir
    $\lambd[fx][f({a}fx)]$ fonksiyonuna değerlenen bir fonksiyondur. Bu ikinci
    fonksiyon kendi başına Church sayısı değildir; ancak $a$ bir Church sayısı
    olduğunda bir Church sayısına indirgenir:}
   (\lambd[a][\lambd[fx][f ({a} f x)]])\,\num{n} & \redone
   \lambd[fx][f ({\num{n}} f x)].
   \intertext{İçteki $\num{n}fx$ terimi bir redextir; çünkü
     $\num{n}$, $\lambd[fx][f^nx]$'tir. Dolayısıyla
     $\num{n}fx\redone f^nx$ ve bütün terim için}
   \fn{Succ}\, \num n & \red \lambd[fx][f(f^n(x))]
\end{align*}
elde edilir; bu da $\num{n+1}$'dir.
\end{proof}

\begin{ex}
  $\fn{Succ}$ fonksiyonunu $\num{0}$, yani $\lambd[fx][x]$ terimine
  uyguladığımızda ne olduğunu görelim. Terimleri bütünüyle açalım:
  \begin{align*}
    \fn{Succ}\, \num{0} & \ident (\lambd[a][\lambd[f][\lambd[x][(f (({a} f) x))]]])(\lambd[f][\lambd[x][x]])\\
    & \redone \lambd[f][\lambd[x][(f (({(\lambd[f][\lambd[x][x]])} f) x))]] \\
    & \redone \lambd[f][\lambd[x][(f ({(\lambd[x][x])} x))]] \\
    & \redone \lambd[f][\lambd[x][(f x)]] \ident \num{1}.
  \end{align*}
\end{ex}

\begin{prob}
  Şu terim
  \begin{align*}
    \fn{Succ}' & \ident \lambd[{n}][\lambd[fx][{n} f (f x)]]
  \end{align*}
  ardıl fonksiyonunu !!{lambda define}s. Nedenini açıklayın.
\end{prob}

\begin{prop}
  \ollabel{prop:add-ld}
  Toplama fonksiyonu $\Add$ !!{lambda definable}.
\end{prop}

\begin{proof}
Şu terimler toplamayı !!{lambda define}:
\begin{align*}
  \fn{Add} & \ident \lambd[{a}{b}][\lambd[fx][{a} f ({b} f x)]]
\intertext{ya da alternatif olarak}
  \fn{Add}' & \ident \lambd[{a}{b}][{a}\, \fn{Succ} \, {b}].
\intertext{İlk tanım şöyle çalışır: $\fn{Add}$ önce iki $a,b$ sayısı alır.
  Sonuç, $f,x$ argümanlarını alıp $af(bfx)$ döndüren bir fonksiyondur.
  $a,b$ sırasıyla $\num n,\num m$ Church sayılarıysa bu,
  $f^{n+m}(x)$'e, yani $f^n(f^m(x))$'e indirgenir. Daha ayrıntılı olarak:}
  (\lambd[{a}{b}][\lambd[fx][{a} f ({b} f x)]])\num n\,\num m & \redone
  \lambd[fx][\num{n}\, f (\num {m}\, f x)] \\
  & \redone \lambd[fx][\num{n}\, f (f^m x)] \\
  & \redone \lambd[fx][f^n (f^m x)] \ident \num {n+m}.
\intertext{İkinci toplama gösterimi $\fn{Add'}$ farklı çalışır. İki Church
  sayısı $\num n,\num m$'ye uygulandığında}
\fn{Add}' \num n \,\num m
& \redone \num{n}\, \fn{Succ}\, \num{m}.
\intertext{Ancak $\num n f x$ her zaman $f^n(x)$'e indirgenir. Dolayısıyla}
  \num{n}\,\fn{Succ}\,\num{m} & \red \fn{Succ}^n(\num{m}).
\end{align*}
$\fn{Succ}$ ardıl fonksiyonunu !!{lambda define}s ve ardıl fonksiyonunu
$m$'ye $n$ kez uygulamak $n+m$ sonucunu verir; dolayısıyla bu terim de
$\num{n+m}$'ye indirgenir.
\end{proof}

\begin{prop}
  \ollabel{prop:mult-ld}
  Çarpma şu terim tarafından !!{lambda definable}:
  \[
  \fn{Mult} \ident \lambd[ab][\lambd[fx][a (b f) x]].
  \]
\end{prop}

\begin{proof}
  $\fn{Mult}$ terimini $\num n,\num m$ Church sayılarına uygulayalım.
  $\fn{Mult}\,\num n\,\num m$, $\lambd[fx][\num n(\num m\,f)x]$'e
  indirgenir. $\num m f$ terimi, $f$'yi argümanına $m$ kez uygulayan bir
  fonksiyonu tanımlar. Dolayısıyla $\num n(\num m f)x$, ``$f$'yi $m$ kez
  uygula'' fonksiyonunu $x$'e $n$ kez uygular. Başka bir deyişle $f$, $x$'e
  $n\cdot m$ kez uygulanır. Ortaya çıkan normal terim $\num{nm}$ Church
  sayısıdır.
\end{proof}

\begin{editorial}
Bu terimi $\eta$-indirgemeyle daha da sadeleştirebiliriz:
\[
  \fn{Mult} \ident \lambd[ab][\lambd[f][a (b f)]].
\]

Ancak bunun için önce $\eta$-indirgemeyi açıklamak gerekir.
\end{editorial}

\begin{prob}
Şu terim çarpmayı !!{lambda define}:
\[
  \fn{Mult}' \ident \lambd[ab][a (\fn{Add}\, a) \num{0}].
\]
Neden çalıştığını açıklayın.
\end{prob}

Üs almanın $\lambd$-terimi olarak tanımı şaşırtıcı ölçüde basittir:
\[
  \fn{Exp} \ident \lambd[be][e b].
\]
İlk $b$ argümanı taban, ikinci $e$ argümanı üstür. Sayı kodlamamız gereği
$ef$, sezgisel olarak $f^e$'dir. Bu tanımı anlamak zor gelirse üs almayı
yinelenen çarpmayla da tanımlayabiliriz:
\[
  \fn{Exp}' \ident \lambd[be][e (\fn{Mult}\, b) \num{1}].
\]

Church sayıları üzerinde öncül ve çıkarma düşündüğümüz kadar basit değildir;
çiftlerin kodlanmasını gerektirir.

